\chapter{Introduction to Fourier Series}
\label{chap:fourier-series}

% ==============================================================================
% SECTION 6.1: OVERVIEW & LEARNING OBJECTIVES
% ==============================================================================
\section{Unit Overview \& Learning Objectives}

Fourier series is one of the most transformative mathematical tools in modern engineering, enabling arbitrary periodic piecewise continuous waveforms to be decomposed into infinite sums of fundamental and harmonic sinusoids. In electrical signal processing, telecommunications, acoustics, mechanical vibration analysis, and heat diffusion PDEs, Fourier analysis converts complex time-domain signals into transparent frequency-domain spectra.

\begin{important}[Core Learning Objectives]
After mastering this unit, the student will be able to:
\begin{enumerate}[leftmargin=1.5em]
    \item \textbf{Analyze Periodic Functions}: Determine fundamental period $T$ and angular frequency $\omega_0 = \frac{2\pi}{T}$.
    \item \textbf{Evaluate Euler Coefficients}: Compute the Fourier coefficients $a_0, a_n, b_n$ over intervals of length $2\pi$ ($(-\pi, \pi)$) and general length $2L$ ($(-L, L)$).
    \item \textbf{Apply Dirichlet's Conditions}: Verify convergence criteria and calculate the sum of the Fourier series at continuous points and jump discontinuities ($\frac{f(x_0^+) + f(x_0^-)}{2}$).
    \item \textbf{Exploit Symmetry}: Expand even functions into Fourier Cosine Series ($b_n = 0$) and odd functions into Fourier Sine Series ($a_0 = a_n = 0$).
    \item \textbf{Construct Half-Range Series}: Derive Half-Range Cosine and Sine Expansions over $(0, L)$.
    \item \textbf{Apply Parseval's Theorem}: Calculate total harmonic power / energy and sum infinite numeric series.
\end{enumerate}
\end{important}

% ==============================================================================
% SECTION 6.2: TRIGONOMETRIC FOURIER SERIES & EULER FORMULAE
% ==============================================================================
\section{Trigonometric Fourier Series \& Euler's Formulae}

\begin{definition}[Fourier Series Representation on $(-\pi, \pi)$]
A periodic function $f(x)$ with period $T = 2\pi$ satisfying Dirichlet's conditions can be represented as:
\begin{equation}
f(x) = \frac{a_0}{2} + \sum_{n=1}^\infty \left[ a_n \cos(nx) + b_n \sin(nx) \right]
\end{equation}
where the \textbf{Euler coefficients} are given by:
\begin{align}
a_0 &= \frac{1}{\pi} \int_{-\pi}^\pi f(x)\,dx \\
a_n &= \frac{1}{\pi} \int_{-\pi}^\pi f(x)\cos(nx)\,dx, \quad n = 1, 2, 3, \dots \\
b_n &= \frac{1}{\pi} \int_{-\pi}^\pi f(x)\sin(nx)\,dx, \quad n = 1, 2, 3, \dots
\end{align}
\end{definition}

\begin{definition}[Fourier Series on Arbitrary Interval $(-L, L)$]
For a periodic function with period $T = 2L$:
\begin{equation}
f(x) = \frac{a_0}{2} + \sum_{n=1}^\infty \left[ a_n \cos\left(\frac{n\pi x}{L}\right) + b_n \sin\left(\frac{n\pi x}{L}\right) \right]
\end{equation}
with Euler coefficients:
\begin{align}
a_0 &= \frac{1}{L} \int_{-L}^L f(x)\,dx \\
a_n &= \frac{1}{L} \int_{-L}^L f(x)\cos\left(\frac{n\pi x}{L}\right)\,dx \\
b_n &= \frac{1}{L} \int_{-L}^L f(x)\sin\left(\frac{n\pi x}{L}\right)\,dx
\end{align}
\end{definition}

% ==============================================================================
% SECTION 6.3: DIRICHLET CONDITIONS & CONVERGENCE
% ==============================================================================
\section{Dirichlet's Conditions \& Pointwise Convergence}

\begin{theorem}[Dirichlet's Conditions for Convergence]
A periodic function $f(x)$ of period $2L$ can be expanded in a convergent Fourier series if on $[-L, L]$:
\begin{enumerate}[leftmargin=1.5em]
    \item $f(x)$ is single-valued and bounded (has at most finite absolute values).
    \item $f(x)$ is piecewise continuous with at most a finite number of jump discontinuities.
    \item $f(x)$ has at most a finite number of local maxima and minima.
\end{enumerate}
\end{theorem}

\begin{theorem}[Value of Fourier Series at Discontinuities]
Let $S(x)$ denote the Fourier series of $f(x)$:
\begin{enumerate}[leftmargin=1.5em]
    \item If $x = x_0$ is a point of \textbf{continuity}, then $S(x_0) = f(x_0)$.
    \item If $x = x_0$ is a point of \textbf{jump discontinuity}, then:
    \begin{equation}
    S(x_0) = \frac{f(x_0^+) + f(x_0^-)}{2}
    \end{equation}
    where $f(x_0^+) = \lim_{h\to 0^+} f(x_0 + h)$ and $f(x_0^-) = \lim_{h\to 0^+} f(x_0 - h)$.
    \item At the interval endpoints $x = \pm L$, the series converges to the periodic boundary average:
    \begin{equation}
    S(\pm L) = \frac{f(-L^+) + f(L^-)}{2}
    \end{equation}
\end{enumerate}
\end{theorem}

% ==============================================================================
% SECTION 6.4: SYMMETRY & HALF-RANGE EXPANSIONS
% ==============================================================================
\section{Symmetry Simplifications \& Half-Range Expansions}

\subsection{Even and Odd Function Expansions}
\begin{itemize}[leftmargin=1.5em]
    \item \textbf{Even Function} ($f(-x) = f(x)$ on $(-L, L)$):
    \begin{equation}
    b_n = 0, \quad a_0 = \frac{2}{L}\int_0^L f(x)\,dx, \quad a_n = \frac{2}{L}\int_0^L f(x)\cos\left(\frac{n\pi x}{L}\right)\,dx
    \end{equation}
    The series contains only cosine terms (\textbf{Fourier Cosine Series}).
    \item \textbf{Odd Function} ($f(-x) = -f(x)$ on $(-L, L)$):
    \begin{equation}
    a_0 = 0, \quad a_n = 0, \quad b_n = \frac{2}{L}\int_0^L f(x)\sin\left(\frac{n\pi x}{L}\right)\,dx
    \end{equation}
    The series contains only sine terms (\textbf{Fourier Sine Series}).
\end{itemize}

\subsection{Half-Range Expansions on $(0, L)$}
When $f(x)$ is defined only on the half-interval $(0, L)$:
\begin{enumerate}[leftmargin=1.5em]
    \item \textbf{Half-Range Cosine Series} (Even Extension):
    \begin{equation}
    f(x) = \frac{a_0}{2} + \sum_{n=1}^\infty a_n \cos\left(\frac{n\pi x}{L}\right), \quad a_0 = \frac{2}{L}\int_0^L f(x)\,dx, \; a_n = \frac{2}{L}\int_0^L f(x)\cos\left(\frac{n\pi x}{L}\right)\,dx
    \end{equation}
    \item \textbf{Half-Range Sine Series} (Odd Extension):
    \begin{equation}
    f(x) = \sum_{n=1}^\infty b_n \sin\left(\frac{n\pi x}{L}\right), \quad b_n = \frac{2}{L}\int_0^L f(x)\sin\left(\frac{n\pi x}{L}\right)\,dx
    \end{equation}
\end{enumerate}

% ==============================================================================
% SECTION 6.5: PARSEVAL'S IDENTITY
% ==============================================================================
\section{Parseval's Identity \& Harmonic Energy}

\begin{theorem}[Parseval's Identity for Fourier Series]
If $f(x)$ has Fourier coefficients $a_0, a_n, b_n$ on $(-L, L)$, then:
\begin{equation}
\frac{1}{L} \int_{-L}^L [f(x)]^2\,dx = \frac{a_0^2}{2} + \sum_{n=1}^\infty (a_n^2 + b_n^2)
\end{equation}
\end{theorem}

\begin{example}[Complete Fourier Series of a Square Wave]
Find the Fourier series for $f(x) = \begin{cases} -k, & -\pi < x < 0 \\ k, & 0 < x < \pi \end{cases}$ with period $2\pi$. Hence deduce $\frac{\pi}{4} = 1 - \frac{1}{3} + \frac{1}{5} - \frac{1}{7} + \cdots$.
\end{example}

\begin{solutionbox}
Since $f(-x) = -f(x)$, $f(x)$ is an \textbf{odd function}.
Therefore, $a_0 = 0$ and $a_n = 0$.
We evaluate $b_n$:
\begin{align}
b_n &= \frac{2}{\pi}\int_0^\pi k \sin(nx)\,dx = \frac{2k}{\pi} \left[ \frac{-\cos(nx)}{n} \right]_0^\pi = \frac{2k}{n\pi}(1 - (-1)^n) \\
&= \begin{cases}
\frac{4k}{n\pi}, & \text{for } n \text{ odd} \\
0, & \text{for } n \text{ even}
\end{cases}
\end{align}
Thus, the Fourier series is:
\begin{equation}
f(x) = \frac{4k}{\pi} \left( \sin x + \frac{\sin 3x}{3} + \frac{\sin 5x}{5} + \frac{\sin 7x}{7} + \cdots \right)
\end{equation}
At $x = \pi/2$, $f(\pi/2) = k$, $\sin(\pi/2) = 1$, $\sin(3\pi/2) = -1$, $\sin(5\pi/2) = 1$:
\begin{equation}
k = \frac{4k}{\pi} \left( 1 - \frac{1}{3} + \frac{1}{5} - \frac{1}{7} + \cdots \right) \implies \frac{\pi}{4} = 1 - \frac{1}{3} + \frac{1}{5} - \frac{1}{7} + \cdots
\end{equation}
\end{solutionbox}

